% Source: tex/chapters/ch05/07_テスト関連測度.tex
\subsubsection{SUT の評価}
SUT の評価では、テスト結果から対象システムの状態を判断する。代表的には、欠陥の種類、欠陥密度、信頼性、運用時の障害傾向、主要業績指標が使われる。

\paragraph{テスト計画と設計に役立つ SUT 測定}
テスト計画には、配備頻度、リードタイム、平均回復時間、変更失敗率などの測度が使われることがある。シフトレフトや DevOps では、これらの測度がテストの計画、管理、結果評価と強く結び付く。

\paragraph{欠陥の種類・分類・統計}
欠陥分類は、どの種類の欠陥が多いか、どの品質属性に関係するかを整理する。利用性欠陥、セキュリティ脆弱性、プライバシー攻撃、サイバーリスクなど、対象に応じて分類体系は異なる。過去の発生頻度は、どこにテスト資源を集中するかを決める根拠になる。

\paragraph{欠陥密度}
欠陥密度は、見つかった欠陥数を SUT の規模で割った値である。規模にはコード行数、機能規模、モジュール数などが使われる。欠陥密度だけで品質を完全には表せないが、比較や傾向把握には役立つ。

\paragraph{寿命試験と信頼性評価}
信頼性評価は、テストをいつ止めるか、次のリリースに進めるかを判断するために使われる。クラウドやフォグ環境では、高い信頼性と可用性を維持することが重要であり、テスト活動もその評価に関わる。

\paragraph{信頼性成長モデル}
信頼性成長モデルは、観測された障害と修正の履歴から、信頼性がどのように向上するかを予測する。多くのモデルでは、障害の原因となる欠陥を修正すれば信頼性が増すと仮定する。障害数に注目するモデルと、障害間隔時間に注目するモデルがある。
